\chapter{Ensemble learning}
\label{ch:ensembles}

\begin{learningobjectives}
After this chapter, you should be able to:
\begin{itemize}
  \item distinguish bagging, random forests, boosting, voting, and stacking;
  \item explain how diversity and dependence govern ensemble benefit;
  \item construct out-of-fold features for leakage-safe stacking; and
  \item compare ensembles using predictive, calibration, and resource evidence.
\end{itemize}
\end{learningobjectives}

\begin{chapterconnection}
A single tree is expressive but unstable. Ensembles turn variation among fitted
learners into a design variable. Some average perturbed learners to reduce
variance; others build a sequence that corrects error. Treating both mechanisms
as ``many models'' misses their different mathematics and failure modes.
\end{chapterconnection}

\section{Averaging helps when errors are not identical}

Let predictors $F_1,\ldots,F_M$ have error variance $\sigma^2$ and pairwise
error correlation $\rho$. Under a simplified equal-correlation model, the
variance of their average is
\[
\rho\sigma^2+\frac{1-\rho}{M}\sigma^2.
\]
Increasing $M$ reduces only the uncorrelated part. Diversity without competence
is not useful, and highly correlated accurate learners may gain little from
averaging.

\section{Bagging and randomized forests}

Bagging fits a base learner to bootstrap samples and averages predictions.
Unpruned trees are natural base learners because their instability creates
variance that averaging can reduce. Out-of-bag observations can provide an
internal diagnostic, though they do not replace an evaluation design matched to
groups, time, or domain shift.

Random forests also sample candidate features at each split. This decorrelates
trees so strong predictors do not force every tree into the same structure.
Extra-trees methods inject more split randomness. The important comparison is
not which name sounds more sophisticated, but how bias, variance, computation,
and stability change on held-out data.

More trees usually stabilize an average rather than overfit in the same way as
more depth. Depth, leaf size, feature sampling, class weighting, and bootstrap
policy still require selection.

\section{Boosting builds an additive model sequentially}

Boosting adds weak learners in stages. AdaBoost increases attention to examples
misclassified by the current ensemble and combines learners with performance-
dependent weights. Gradient boosting views the model as an additive expansion
and fits each new learner toward the negative gradient of the loss.

For squared error, the negative gradient is proportional to the residual, which
makes ``fit the residuals'' exact. For other differentiable losses, pseudo-
residuals replace ordinary residuals. Learning rate and number of stages trade
step size against sequence length; tree depth controls interaction order.

Histogram-based implementations make strong tabular learners practical by
binning candidate thresholds and optimizing computation. Libraries differ in
category handling, missingness, regularization, determinism, and serialization.
An external package is a dependency with a contract, not merely an import.

\conceptcheck{Why can adding more trees to a random forest and adding more
stages to a boosted ensemble have different overfitting behavior?}

\section{Voting and stacking combine heterogeneous learners}

Hard voting combines predicted labels and discards confidence. Soft voting
averages probabilities, which requires compatible class order and meaningful
calibration. Weight selection is another adaptive choice and belongs inside
validation.

Stacking trains a meta-learner on predictions from base learners. If base
predictions for a training row come from models fitted on that same row, the
meta-learner sees optimistically biased features. Correct stacking constructs
out-of-fold predictions: each training row is predicted by a base model that did
not train on it. Base learners are then refit on the full training data for
future inference.

\begin{warningbox}
The stacking split must preserve the original dependence structure. Ordinary
row-wise folds do not become safe merely because they are used inside an
ensemble. Groups, time, and preprocessing boundaries still apply at every
level.
\end{warningbox}

\section{Importance and explanation become more delicate}

An ensemble can improve prediction while making one global explanation harder.
Impurity importance inherits the biases of tree splits. Permutation importance
measures performance loss under a specified perturbation and can understate
redundant correlated features. Local additive explanation methods approximate
or allocate behavior under background-distribution assumptions.

Use interpretation as a scoped diagnostic. It cannot promote an associative
ensemble into a causal mechanism. Stability across folds, batches, and plausible
perturbations is evidence worth reporting.

\begin{misconceptionbox}
\textbf{Misconception:} ensembles always outperform their members. 
\textbf{Correction:} benefit depends on base competence, error dependence,
selection, and the deployment distribution. Combination can preserve shared
bias, amplify leakage, damage calibration, and exceed resource constraints.
\end{misconceptionbox}

\section{Champion-challenger evaluation}

The champion is the approved artifact currently serving the contract. A
challenger is evaluated under the same split, metrics, calibration procedure,
and resource budget. Promotion should require a practically meaningful gain,
not merely the largest point estimate among many trials.

Measure serialized size, fit cost, batch throughput, tail latency, memory, and
energy where material. Verify deterministic behavior under the supported
runtime or document unavoidable variation. A smaller calibrated model may be a
better system component than a marginally more accurate ensemble.

\begin{workedexample}{forest and boosting challengers}
For future-batch degradation prediction, the logistic champion is compared with
a random forest and histogram gradient booster. All candidates receive the same
training-only preprocessing and outer temporal evaluation. Hyperparameters are
chosen within earlier batches.

The booster improves average precision, but its raw probabilities overstate
event frequency. Calibration is fitted without touching the final batch. The
release report compares the calibrated system with the champion on decision
cost, batch-level uncertainty, model size, and prediction latency. Promotion is
withheld when the apparent gain disappears on the newest batch. The result is a
successful evaluation, not a failed project.
\end{workedexample}

\section{Exercises}
\begin{enumerate}
  \item Derive the equal-correlation variance formula for an average of $M$
  predictors.
  \item Draw the data flow required to generate leakage-safe stacking features.
  \item Compare random forest and gradient boosting when base trees are made
  deeper. Explain the different mechanisms.
  \item Write a promotion gate that includes predictive quality, calibration,
  artifact size, latency, and subgroup regression limits.
\end{enumerate}

\begin{chapterreview}
Ensembles improve learning through structured diversity. Bagging and randomized
forests reduce variance by averaging perturbed learners. Boosting constructs an
additive sequence guided by loss. Voting and stacking combine heterogeneous
models, but stacking requires out-of-fold features. Accuracy gains must survive
calibration, dependence-aware evaluation, interpretation limits, and resource
budgets.
\end{chapterreview}

\begin{notebooklink}
Lecture 14 notebook 00 maps forests, boosting, voting, and stacking. Companion
laboratories build a leakage-safe benchmark and a champion-challenger report.
\end{notebooklink}
